\documentclass[a4paper,12pt]{article}

\usepackage[utf8]{inputenc}
\usepackage[T1]{fontenc}
\usepackage[italian]{babel}
\usepackage{geometry}
\usepackage{amsmath}
\usepackage{booktabs}
\usepackage{xcolor}
\usepackage{listings}
\usepackage{enumitem}

\geometry{margin=2.5cm}

\lstset{
    basicstyle=\ttfamily\small,
    backgroundcolor=\color{gray!10},
    frame=single,
    breaklines=true,
    columns=fullflexible
}

\title{Scaling del problema QUBO} 
\author{}
\date{May 2024}

\begin{document}

\maketitle
\newpage

\section{Scaling del problema QUBO}
Come si è visto una QUBO ha forma:
\[
f(x) = x^T Q x
\]

dove \(Q\) contiene i coefficienti lineari e quadratici del problema.  
Se moltiplichiamo tutta la matrice per una costante positiva \(\alpha\):

\[
Q' = \alpha Q
\]

allora la nuova funzione diventa:

\[
f'(x) = x^T Q' x
\]

Sostituendo \(Q' = \alpha Q\), otteniamo:

\[
f'(x) = x^T (\alpha Q) x
\]

Da cui:

\[
f'(x) = \alpha x^T Q x = \alpha f(x)
\]

Nei test si è usata: 
\[
\alpha \in [0.1, 0.2]
\]
\[\]
Per esempio, se due configurazioni hanno energie:

\[
E(x_1) = 100, \qquad E(x_2) = 110
\]

la differenza energetica è:

\[
\Delta E = E(x_2) - E(x_1) = 110 - 100 = 10
\]

Se scaliamo tutto con \(\alpha = 0.1\), otteniamo:

\[
E'(x_1) = \alpha E(x_1) = 0.1 \cdot 100 = 10
\]

\[
E'(x_2) = \alpha E(x_2) = 0.1 \cdot 110 = 11
\]

e quindi:

\[
\Delta E' = E'(x_2) - E'(x_1) = 11 - 10 = 1
\]

La soluzione migliore è ancora \(x_1\), perché continua ad avere energia minore 
rispetto a \(x_2\). Tuttavia, il salto energetico tra \(x_1\) e \(x_2\) è diventato 
molto più piccolo.

\section{Perche scalare il problema?}
Nel simulated annealing, la probabilità di accettare una mossa peggiorativa è data da:
\[
P = e^{-\frac{\Delta E}{T}}
\]

Quindi, a parità di temperatura T, se $\Delta E$ è grande, la probabilità diventa molto piccola. 
Se invece lo scaling riduce $\Delta E$, la probabilità aumenta.
\\Esempio con $T$=1:
\[
e^{-10} \approx 0.000045, \qquad e^{-1} \approx 0.367
\]

Si deduce che dopo lo scaling, una mossa leggermente peggiorativa può ancora essere accettata 
con probabilità significativa. Questo permette all’annealer di esplorare meglio lo 
spazio delle soluzioni, invece di bloccarsi subito nel primo minimo locale trovato.

Tuttavia, uno scaling eccessivo può rendere il paesaggio energetico troppo piatto, riducendo la 
capacità dell'annealer di distinguere soluzioni buone da soluzioni peggiori.

\end{document}